\documentclass{article}
\usepackage[utf8]{inputenc}
\usepackage[T1]{fontenc}
\usepackage{geometry}
\usepackage{graphicx}
\usepackage{float}
\usepackage{hyperref}
\usepackage[english,indonesian]{babel}

\title{Optimasi Pemrosesan Template LaTeX Menggunakan Pendekatan Berbasis AI}
\author{Academic Support Agent}
\date{\today}

\begin{document}
\maketitle

\begin{abstract}

\end{abstract}

\chapter{PENDAHULUAN}


\chapter{TINJAUAN PUSTAKA}


% Original content (commented out):
%  \documentclass[a4paper, twoside]{article}

\usepackage{amsmath, amsthm, bm, amssymb}
\usepackage{mathtools}
\usepackage{graphicx}
\usepackage{xcolor}
\usepackage[utf8]{inputenc}
\usepackage{authblk}
\usepackage{titlesec} % personalisar as seções
\usepackage{amsmath,,amsthm} % pacotes matemáticos básicos
\usepackage{indentfirst} % indenta sempre o primeiro parágrafo
\usepackage{geometry} % Margens
\usepackage{icomma} % corrige espaçamento quando vírgula é utilizada como separador decimal
\usepackage{booktabs} % melhorar exemplo de tabela
\usepackage{tabularx} % tabelas
\usepackage{multirow}
\usepackage{fancyhdr} % Cabeçalhos
\usepackage{natbib} % \citep, \citeyear, etc
\usepackage{caption} % para definir a fonte e alinhamento das legendas
\usepackage[brazil]{babel}
\usepackage{hyperref}
\usepackage{etoolbox}
\usepackage{enumitem}
\usepackage{float}

%%%%%%%%%%%%%%%%%%%%%%%%%%%
% Configurações da página %
%%%%%%%%%%%%%%%%%%%%%%%%%%%
\newgeometry{
  bottom = 2.5cm,
  top = 2.5cm,
  left = 3.0cm,
  right = 3.0cm,
  headheight = 27pt
}
\linespread{1.5} % equivale ao entrelinhas de 1.5 do word!
\setlength{\parindent}{1.0cm} % Tamanho da indentação do parágrafo.
\AtBeginEnvironment{thebibliography}{
\linespread{1}
\selectfont
\setlength{\bibsep}{12pt}
}
\setlist{nosep}
\setlength{\headheight}{27.31668pt}
%%%%%%%%%%%%%%%%%%%%%%%%%%%%%%%%%%%%%%%%%%
% Configurações dos cabeçalhos e rodapés %
%%%%%%%%%%%%%%%%%%%%%%%%%%%%%%%%%%%%%%%%%%
\pagestyle{fancy}
\fancyhf{}

% Header da primeira página.
\fancypagestyle{firstpage}{%
 \fancyhf{}%
 \fancyhead[L]{\footnotesize{\emph{IV JMMA \\ Jornada de Matemática e Matemática Aplicada \\ Santa Maria - RS, 15, 16 e 17 de Novembro de 2023}}}%
 \fancyfoot[R]{\thepage}
}

%Header das páginas subsequentes.
\fancyhf{}
\renewcommand{\headrulewidth}{0pt}
\renewcommand{\footrulewidth}{0pt}
\fancyhead[CE]{{ \scriptsize IV JMMA, Santa Maria - RS, 15, 16 e 17 de Novembro de 2023}}
\fancyhead[CO]{{ \scriptsize IV JMMA, Santa Maria - RS, 15, 16 e 17 de Novembro de 2023}}
\rfoot{\thepage}

%%%%%%%%%%%%%%%%%%%%%%%%%%%%%%%%%%%%%%%%%%%%%%%%%%%%%%%%%%%%%%%%%%%%%%%%%%5%%%%%%%%%%%%%%%%%%%%%%%%%%%%%%%%%%%%%%
%   NÃO MODIFIQUE NENHUMA PARAMETRO ACIMA, APENAS PREENCHA O TÍTULO, OS AUTORES E AS AFILIAÇÕES ABAIXO          %
%%%%%%%%%%%%%%%%%%%%%%%%%%%%%%%%%%%%%%%%%%%%%%%%%%%%%%%%%%%%%%%%%%%%%%%%%%5%%%%%%%%%%%%%%%%%%%%%%%%%%%%%%%%%%%%%%

\title{IV Jornada de Matemática e Matemática Aplicada - Instruções para autores}

\author[1]{Autor 1 - Nome e sobrenome}
\author[2]{Autor 2 - Nome e sobrenome}
\affil[1]{Afiliação autor 1}
\affil[2]{Afiliação autor 2}

\date{}

\begin{document}

\maketitle
\thispagestyle{firstpage}
    
\renewcommand{\abstractname}{Resumo}
\begin{abstract}
Escrever resumo em português
\end{abstract}

\section{Introdução}

%Serão aceitos trabalhos para apresentação oral (resumos de 5 a 7 páginas) nas sessões técnicas, voltados aos alunos de pós-graduação, com resultados parciais ou finais, além de trabalhos para apresentação na forma de pôsters (resumo de 1 a 2 páginas), voltados aos alunos de graduação.

Escreva seu texto aqui

\newpage

%%%%%%%%%%%%%%%%%%%%%%%%%%%%%%%%%%%%%%%%%%%%%%%%%%%%%%%%%%%%%%%%%%%%%%%%%%5%%%%%%%%%%%%%%%%%%%%%%%%%%%%%%%%%%%%%%
%%%%%%%%%%%%%%%%%%%%%%%%%%%%%%%%%%%%%%%%%%%%%%%%%%%%%%%%%%%%%%%%%%%%%%%%%%5%%%%%%%%%%%%%%%%%%%%%%%%%%%%%%%%%%%%%%

\renewcommand{\refname}{Referências Bibliográficas}
\bibliographystyle{acm}

%Utilize a bibliografia com o "thebibliography" ou usando o bibfile
% \begin{thebibliography}{0}
% \bibitem{}
% \end{thebibliography}



\end{document}


\end{document}